\section{引子}

上一章我们只抽象了一个事件。
本章我们将抽象两个事件。
先考虑一个实际例子。

~

\begin{example}
设甲、乙两人约定在1:00～2:00之间在某处碰头，若有一个先到，等10分钟后如无人则离去，设两人会在1:00～2:00之间任意时刻抵达约定地点，计算两人碰面的概率。
\end{example}

解：

我们从直观感觉出发，首先可以确定，甲、乙两人各自出现在约定地点的概率在1:00～2:00间均匀分布，从图形来看，即下图阴影区的占比。
这是一个典型的几何概率，可得：
\begin{figure}[h]
\centering
\includegraphics[height=2.4cm]{3.1.png}
\end{figure}
\[
P=1-\frac{2\cdot \left( \frac{1}{2}\cdot 50\cdot 50 \right)}{60\cdot 60}=\frac{11}{36}
\]
虽然结果是对的，但缺乏严密地推理，也不利于推广到任意情况，如甲若偏好来得迟，概率结果会如何等等。

~

本章就是将两个变量的情况加以详细讨论。




